\section{Counting}
\label{sec:counting}

\subsection{Pipeline}

Detections from YOLO11m at 640\,px are associated by BoT-SORT~\cite{aharon2022botsort} with
sparse-optical-flow global motion compensation --- mandatory here, since the whole frame
translates as the vehicle moves --- producing candidate tracks. A confirmation stage then
decides which candidates are real, distinct animals. The count is the number of accepted
candidates.

The confirmation baseline is a three-parameter rule: accept a candidate if it persists for
at least $m$ frames, spans at least $s$ seconds, and has top-$k$ mean detector confidence at
least $c$. Under the held-out protocol (\cref{sec:heldout}) the swept values are
$m\ge20$, $s\ge0$, $c\ge0.65$.

We additionally recover \emph{orphan} detections --- those the tracker returns without a
track identity --- by linking detections within a temporal gap and a scale-normalized
distance into pseudo-tracks. Without this, single-frame animals are discarded entirely; with
it, 17 animals are recovered.

\subsection{The headline result}

\begin{table}[t]
\centering\small
\setlength{\tabcolsep}{4pt}
\begin{tabular}{lccc}
\toprule
Scope & MAE & Bias & Counted \\
\midrule
Fit videos (detector saw these) & 1.53 & $-0.05$ & 99.3\% \\
\textbf{Held out (13 unseen videos)} & \textbf{2.38} & $\mathbf{-1.92}$ & \textbf{55/83 = 66.3\%} \\
All 32 videos (optimistic) & 1.88 & $-0.81$ & 89.0\% \\
\bottomrule
\end{tabular}
\caption{\textbf{End-to-end counting.} The held-out row is the honest number: the rule is
swept on the 19 detector-training videos, frozen, and applied to the 13 unseen ones. Bias is
consistently negative --- the system under-counts and essentially never over-counts, which is
the conservative direction for a wildlife survey.}
\label{tab:counting}
\end{table}

\Cref{tab:counting} is the end-to-end result: \textbf{MAE 2.38 animals per video and 66.3\%
of unseen animals counted}, with a bias of $-1.92$. The system under-counts systematically
and over-counts almost never, which for management purposes is the safe direction --- a
survey that never inflates its estimate is preferable to one that occasionally does.

Per site, the error concentrates in dense groups: held-out MAE is 4.33 at SHB (the
group-heavy site), 2.20 at TON, 2.67 at SHW, and 1.00 at MAS.

Under the stricter identity-matched criterion of \cref{sec:decomposition} --- animals
covered by an \emph{accepted} track, rather than per-video count agreement --- the figure is
47/83 = 56.6\%. The 10-point gap is coincidental agreement: accepted tracks that do not land
on an animal offsetting animals that were missed. We report the capped aggregate as our
headline for comparability with the counting literature, and the identity-matched figure so
the difference is on record.

\subsection{Improving candidate generation makes counting worse}
\label{sec:inversion}

This is the paper's central experimental result. We improved the candidate pool four
independent ways and measured the effect on the decomposition and on the final count.

\begin{table*}[t]
\centering\small
\begin{tabular}{llrccccc}
\toprule
Pool & Intervention & Candidates & \reachedq & \primaryq & \countedq & MAE & Ceiling \\
\midrule
C & Orphan recovery, conf 0.10 (operating point) & 7,008 & 74/83 & 66/83 & \textbf{55/83 = 66.3\%} & \textbf{2.38} & 79.5\% \\
E & \quad + track-initialization gate removed, conf 0.02 & 18,349 & 79/83 & 69/83 & 54/83 = 65.1\% & 2.54 & 83.1\% \\
F & \quad + NMS IoU $0.50\!\to\!0.90$ & 90,239 & 79/83 & 70/83 & 41/83 = 49.4\% & 3.46 & 84.3\% \\
G & \quad + appearance re-identification & 27,679 & \textbf{82/83 = 98.8\%} & \textbf{73/83 = 88.0\%} & 52/83 = 62.7\% & 2.69 & \textbf{88.0\%} \\
\bottomrule
\end{tabular}
\caption{\textbf{Every improvement to candidate generation degrades the count.} Pool G
reaches 82 of 83 unseen animals and gives 73 their own track --- the best pool we built, at
less than a third of pool F's candidate count --- yet counts fewer animals than pool C, which
reaches only 74. ``Ceiling'' is \primaryq/83, the maximum a one-count-per-track confirmation
stage could achieve.}
\label{tab:pools}
\end{table*}

\begin{figure}[t]
  \centering
  \includegraphics[width=\linewidth]{figures/fig2_inversion.pdf}
  \caption{\textbf{The inversion.} As candidate generation improves, more animals become
  reachable and fewer are counted. Both series are animals out of 83 and share one axis.}
  \label{fig:inversion}
\end{figure}

\Cref{tab:pools} and \cref{fig:inversion} give the result. Each intervention does what it
was designed to do:

\begin{itemize}\itemsep2pt \parskip0pt
\item \textbf{Lowering detector confidence} 0.10$\to$0.02 and \textbf{removing the tracker's
track-initialization gate} raised \reachedq from 74 to 79. We note in passing that lowering
confidence \emph{alone} recovered zero additional animals: the gate, not the threshold, was
the binding constraint --- BoT-SORT will extend a track from a low-confidence detection but
will not \emph{start} one, so faint animals were invisible to counting regardless of
detector threshold.
\item \textbf{Loosening NMS} to IoU 0.90 addressed animals standing close enough that one
box suppressed another; it added 4 \primaryq at the cost of $5\times$ the candidates.
\item \textbf{Appearance re-identification} produced the best pool by every measure: 82 of
83 animals reached, 73 with their own track, at 27,679 candidates --- fewer than a third of
pool F's for four more \primaryq. It genuinely works: it merges fragments of one animal
rather than merely adding candidates.
\end{itemize}

And the count falls, monotonically in candidate volume: 55, 54, 52, 41. The confirmation
stage cannot exploit a pool it did not shrink. Each intervention adds far more negatives
than positives --- pool G has 219 positive tracks against 27,460 negatives --- and a
three-parameter rule tuned to survive that ratio must be aggressive enough to reject real
animals along with the noise.

\subsection{Confirmation is a supervision problem, not a capacity problem}
\label{sec:learned}

If the rule is the bottleneck, replacing it should help. We tested that under three
protocols and it does not.

\begin{table}[t]
\centering\small
\begin{tabular}{llcc}
\toprule
Pool & Confirmer & MAE & Counted \\
\midrule
\multirow{4}{*}{C} & \textbf{Hand-tuned rule (3 params)} & \textbf{2.38} & \textbf{55/83} \\
 & Gradient boosting & 2.85 & --- \\
 & Logistic regression & 3.00 & --- \\
 & Temporal transformer ($\sim$60k params) & 2.85 & 49/83 \\
\midrule
\multirow{4}{*}{G} & Hand-tuned rule & 2.77 & 50/83 \\
 & Gradient boosting & 2.92 & --- \\
 & Logistic regression & 3.15 & --- \\
 & Temporal transformer & 2.77 & 51/83 \\
\bottomrule
\end{tabular}
\caption{\textbf{Held-out confirmation comparison.} All methods trained and thresholded on
the 19 detector-training videos, reported on the 13 unseen. Learned confirmers lose on pool
C and tie on pool G. Logistic regression on pool G fails in the opposite direction,
predicting 102 animals where 83 exist.}
\label{tab:confirmers}
\end{table}

\Cref{tab:confirmers} gives the held-out comparison. On pool C every learned method loses;
on pool G the transformer ties. Notably, every learned variant \emph{under-counts harder}
than the rule, so they are rejecting more candidates rather than recovering the animals the
rule rejects.

Two further protocols agree. Under 8-fold cross-validation over videos
(\cref{tab:capacity}) the rule again wins, and the pattern in that table is diagnostic:
performance degrades monotonically with model capacity, and shrinking a gradient-boosted
model to depth-1 stumps ($\sim$40 effective parameters) recovers nearly the whole gap and
achieves the best RMSE of any method. This is a variance limit, not an architectural one.

\begin{table}[t]
\centering\small
\begin{tabular}{lrcc}
\toprule
Method & Params & MAE & RMSE \\
\midrule
\textbf{Hand-tuned rule} & 3 & \textbf{1.88} & 2.99 \\
Capacity-matched stumps & $\sim$40 & 1.97 & \textbf{2.78} \\
Gradient boosting (depth 3) & $\sim$10$^3$ & 2.41--2.66 & 3.60 \\
Temporal transformer & $\sim$60k & 2.81 & 4.86 \\
Logistic regression & 13 & 3.41 & 4.86 \\
Learned accept + clustering & $\sim$10$^3$ & 3.66 & 5.99 \\
Soft-count MLP & $\sim$5k & 3.89--4.44 & 6.34 \\
\bottomrule
\end{tabular}
\caption{\textbf{Capacity, not architecture, decides the confirmation stage} (8-fold CV over
videos). Error grows monotonically with parameter count; capacity-matched stumps recover
nearly the whole gap and give the best RMSE, i.e.\ fewer catastrophic per-video errors.}
\label{tab:capacity}
\end{table}

We also trained directly on the counting objective, letting fragments share probability mass
so the loss is $|\sum_i p_i - c|$ rather than per-track cross-entropy. It was the worst
variant tested, which rules out ``the rule wins because it optimizes MAE directly'' as the
explanation.

\subsection{Calibrated per-animal confidence}

A survey system that cannot say how sure it is cannot be audited. We calibrate a per-track
posterior for ``this track is on a real animal'' and report it with every counted animal:
\textbf{92.1\% of counted animals carry confidence $\ge0.80$} (mean 0.965, AUC 0.893), and
793 low-confidence tracks are flagged for human review rather than auto-counted. We also
propagate per-track probabilities into a Poisson-binomial standard deviation on each video's
count, giving an uncertainty interval rather than a point estimate.

Raw detector confidence cannot serve this purpose: no track in the corpus exceeds 0.90 and
the median counted track sits at 0.74, because the animals are 27 pixels. The requirement is
only meaningful against a calibrated posterior.

\begin{figure}[t]
  \centering
  \includegraphics[width=0.92\linewidth]{figures/fig3_per_deer_evidence_sheet.jpg}
  \caption{\textbf{Per-animal evidence.} Every counted animal is exported as a contact sheet
  of its own track with timestamps and per-frame confidence, so a human auditor can verify
  any individual count. Confidence $\approx$0.82 for this track.}
  \label{fig:evidence}
\end{figure}
